\section[Dilemme de l'infrastructure]{Le dilemme de l'infrastructure: entre le sur-mesure et le standard}
La refonte de la bibliothèque numérique de la \ac{BIU Cujas} soulève une interrogation stratégique dépassant la simple technique informatique : quel modèle d'infrastructure privilégier pour garantir à la fois l'indépendance et l'adéquation aux besoins scientifiques et la pérennité de l'outil ? En ingénierie documentaire, c'est une question classique opposant la logique du développement spécifique, incarné par le prototype Codex, à celle de l'adoption d'outil standard, comme Omeka S.

\subsection{L'analyse du sur-mesure}
Le choix initial de la Bibliothèque Cujas, en développant le prototype Codex, offre une flexibilité totale et une forte adaptation aux besoins métiers. Cela évite également la dépendance finiancière et technique vis-à-vis d'un prestataire externe, \enquote{qui facture cher la moindre évolution}. Conçu en Python-Django, Codex offre l'avantage de s'interfacer dans la façon de travailler des experts métiers. Ce couplage permettrait  d'automatiser certaines tâches de la chaîne de numérisation de Cujas, comme la génération de métadonnées \ac{METS} ou \ac{XML}, ou des transferts sécurisés vers le serveur d'archivage du \ac{CINES}.

Si l'adaptation aux besoins est forte, ce mode de développement démontre une extrême vulnérabilité humaine, identifiée, par le directeur de la bibliothèque. Il le perçoit comme un risque de pérennité de l'infrastructure en bâtissant la future bibliothèque numérique uniquement sur les compétences d'une seule personne. En effet, Olesea Dubois confirme cette analyse : \enquote{le plus gros soucis avec les outils maisons c'est qu'ils sont très souvent associés à des personnes et à des procédures qui leur sont propres}\footannexe{\nameref{annexe_A_rapport_audit}}[pour voir l'entretien réalisé auprès d'Olesea Dubois]. Un autre inconvénient mis en avant par la responsable de la bibliothèque numérique de Sciences Po est que \enquote{généralement, il y a peu ou pas assez de documentations entraînant ainsi une durée de vie assez courte de la bibliothèque numérique} \footannexe{\nameref{annexe_A_rapport_audit} pour voir l'entretien avec Olesea Dubois.}. Effectivement, le peu de documentation disponible entraîne un risque pour la pérennité de l'outil, et empêche le maintien sur le long terme de la bibliothèque numérique. Avoir recours à un logiciel sur-mesure entraîne une difficulté de recrutement, car l'institution se coupe d'un vivier de compétence, rendant rapidement le code sur-mesure inexploitable en cas du départ de son concepteur, condamnant le système à une obsolescence rapide. La crainte de la plupart des institutions patrimoniales y compris la bibliothèque Cujas est de reproduire la situation qui a entraîné la disparition de la bibliothèque numérique Medica. C'est pour pallier à cela qu'elles ont choisi de se tourner vers des outils standardisés afin de disposer d'un vivier de personnel.

La tension entre le développement d’un prototype sur-mesure et l’adoption d’un standard communautaire comme Omeka S s'inscrit au cœur des réflexions sur la souveraineté technique et la maîtrise de l'obsolescence logicielle au sein des institutions patrimoniales\footcite{bermesWebSemantiqueBibliotheque2013}. D'un côté, le développement interne d’une solution ad hoc, à l'instar du prototype \textit{Codex} conçu sous Python-Django, s'impose comme une réponse souveraine et immédiate face à l'urgence fonctionnelle : il permet de pallier la fragilité de CujasNum sous l'infrastructure obsolète \ac{XTF} et d'assurer une continuité absolue du service public de la recherche face aux risques de sécurité informatique. 

Toutefois, cette autonomie technique se heurte rapidement aux limites de la maintenance à long terme d'un outil isolé. À l'ère de la Transition bibliographique, la rapidité des évolutions technologiques rend en effet indispensable l'adossement à une communauté active pour éviter qu'une institution ne se retrouve marginalisée par des services rapidement dépassés. C'est ici que l'adoption d'un standard libre et communautaire tel qu'Omeka S révèle sa force stratégique. En s'appuyant sur un code source ouvert, l'établissement lutte activement contre l'obsolescence logicielle souvent entretenue par les éditeurs de logiciels propriétaires\footcite{creusotQuellesPerspectivesPour}. De plus, en permettant un déploiement et un hébergement locaux — contrairement aux offres propriétaires en mode SaaS (\textit{Software as a Service}) sur le cloud —, un outil comme Omeka S garantit à l’institution la maîtrise absolue de ses données et de ses choix d'archivage, réalisant un compromis idéal entre les impératifs du modèle sémantique sémantisé (RDF, Dublin Core) et l'utilisabilité de l'interface\footcite{chaineuxLentreeDansWeb, creusotQuellesPerspectivesPour}.


\subsection{L'analyse du standard communautaire}
Face au risque d'isolement technique du sur-mesure, les retours d'expérience des responsables de bibliothèque numérique des institutions patrimoniales comme la Bibliothèque Sainte-Geneviève, la Bibliothèque de Sciences Po ou l'\ac{ENS}, démontrent un abandon progressif des outils sur-mesure au profit d'Omeka S. Financé par des subventions de recherche, de fondations et soutenu par une communauté, Omeka S offre des garanties majeures de durabilité. Comme le résume Thomas Chaineux, ce logiciel représente \enquote{un compromis entre les impératifs du modèle et l'ergonomie}, s'imposant pour les \ac{BU} moyennes comme \enquote{une porte d'entrée bienvenue} grâce à sa gratuité et la force de son réseau communautaire \footcite{chaineuxLentreeDansWeb}. Contrairement à Django qui nécessite d'avoir des compétences informatiques avancées pour développer le logiciel et l'intégrer dans le web sémantique. La particularité d'Omeka S est d'avoir été pensé pour s'intégrer directerement ou plus facilement dans le web de données. Il s'appuie sur une \enquote{modélisation des métadonnées sous \ac{RDF}} combinées à des vocabulaires standardisés, comme le Dublin Core, qui favorise \enquote{une meilleure exposition des données}. Il facilite l'exposition des métadonnées, leur exportation et leur interopérabilité. 
 
De plus, un autre des avantages pour les catalogueurs est de \enquote{bénéficier d'un back-office documenté, collaboratif et modulaire}. Toutefois, malgré les avantages certains, l'environnement Omeka S s'accompagne de contraintes réelles. Bien qu'open-source, le logiciel est \enquote{gratuit}, mais il exige des ressources pour fonctionner et évoluer. Omeka S demande des ressources internes pour en assurer l'hébergement, le paramétrage et la maintenance technique. Gilles Clément, l'un des concepteurs de NuBIS, témoigne de cette limite : \enquote{À l'usage, chaque nouvelle version d'Omeka S cassait quelque chose[...]. Nous étions confrontés à beaucoup de bidouillages en raison de l'instabilité des modules tiers développés par des prestataires différents, ce qui rendait le code lourd et difficile à maintenir} \footannexe{\nameref{annexe_A_rapport_audit} pour voir l'entretien avec Cécile Obligi et Gilles Clément de la \ac{BIS}.}. Cette instabilité fait écho à l'analyse de Thomas Chaineux, qui situe Omeka S dans \enquote{une position hybride dans le Web de données}. Faute d'émettre des \ac{URI} nativement pérennes à l'échelle globale sans développements tiers, le logiciel génère par défaut des \enquote{noeuds blancs} dont la portée reste limitée au contexte local \footcite{chaineuxLentreeDansWeb}. 
 
Face à cet obstacle, la solution mise en place par la \ac{BIS} pour le résoudre, a été de séparer techniquement le front du back, \enquote{rendant ainsi le code-source plus maintenable dans le temps} \footannexe{\nameref{annexe_A_rapport_audit} pour voir l'entretien avec Cécile Obligi et Gilles Clément de la \ac{BIS}.}. En effet, à chaque mise à jour, le module complémentaire déréglait les paramétrages du front-office. La \ac{BIS} a donc découplé son architecture en décidant de \enquote{développer leur back-office sous Omeka S et du front sous Astro.js} \footannexe{\nameref{annexe_A_rapport_audit} pour voir l'entretien avec Cécile Obligi et Gilles Clément de la \ac{BIS}.}. Cette distinction rend étanche les couches logicielles permettant ainsi le bon fonctionnement de NuBIS et la satisfaction des équipes qui en sont en charge. Malgré cet obstacle rencontré par certaines institutions, l'open-source, contrairement au sur-mesure, a cette \enquote{notion de mutualisation rendant plus facile l'idée du maintien de cette structure} \footannexe{\nameref{annexe_A_rapport_audit} pour voir l'entretien avec Vincent Tardif de l'\ac{ENS}.}. Le choix de la \ac{BIS} du framework Astro.js répond à des exigences d'éco-conception et d'accessibilité numérique, évitant ainsi l'interrogation de la base à chaque clic de l'usager. L'avantage majeur de ce type de fonctionnement est que les montées de versions d'Omeka S n'impactent plus l'affichage public, sécurisant mieux le back-office le rendant plus maintenable dans le temps et pérenne.